\chapter{Background and Theoretical Foundations}
\label{ch:background}

This chapter establishes the technical vocabulary and the theoretical results on
which the rest of the thesis depends. It moves from the general (blockchains and
smart contracts) to the specific primitives the system composes (authenticated
encryption, elliptic-curve and threshold cryptography, zero-knowledge proofs,
stealth addresses, and permanent storage), and finishes with the analysis of
trusted execution environments that motivates the project's central design
decision. A reader already fluent in applied cryptography and Ethereum
development may skim this chapter and refer back to it as needed.

\section{Blockchains, Ethereum, and Smart Contracts}
\label{sec:bg-blockchain}

A \emph{blockchain} is an append-only, replicated ledger maintained by a network
of mutually distrusting nodes that reach agreement on its contents through a
consensus protocol. Its defining property for this work is not its data
structure but its \emph{trust model}: state transitions are validated
independently by every participant against public rules, so that the correctness
of the ledger does not depend on trusting any single operator. This is precisely
the property a censorship-resistant system needs.

Ethereum~\cite{buterin2014ethereum, wood2014ethereum} generalised the ledger
into a \emph{world computer}. In addition to recording balances, Ethereum nodes
run the \emph{Ethereum Virtual Machine} (EVM), a deterministic state machine that
executes programs called \emph{smart contracts}. A smart contract is code,
typically written in Solidity~\cite{solidity}, that is deployed to an address and
thereafter executes exactly as written whenever a transaction calls it. Once
deployed, a contract's logic is immutable (absent an explicit upgrade mechanism)
and its execution is verified by the whole network. Contracts can hold and
transfer the native currency (ether), maintain persistent storage, and emit
\emph{events}---append-only log records that off-chain applications can index.

Two characteristics of the EVM are central to this thesis. First, every
operation costs \emph{gas}, a fee paid in ether that bounds computation and
prevents denial-of-service; an efficient contract is therefore a cheaper
contract, which makes minimising on-chain data a concrete engineering goal.
Second, contract state is \emph{public}: anything written to storage or emitted
as an event is visible to everyone. A privacy-preserving application must
therefore be designed so that nothing sensitive is ever placed on-chain---only
commitments, hashes, and references to data held elsewhere. This principle,
``no plaintext on-chain'', recurs throughout the design.

An \emph{externally owned account} (EOA) is controlled by a private key and is
how human users transact. The Ethereum ecosystem has standardised many
behaviours through \emph{Ethereum Requests for Comment} (ERCs); this work uses
ERC-5564 (stealth addresses), discussed in Section~\ref{sec:bg-stealth}.

\section{Symmetric Encryption and Authenticated Encryption}
\label{sec:bg-aead}

Confidentiality of the source's evidence rests on \emph{symmetric encryption}, in
which the same secret key both encrypts and decrypts data. The system uses the
\emph{Advanced Encryption Standard} (AES)~\cite{nistAES} with a 256-bit key,
operating in \emph{Galois/Counter Mode} (GCM)~\cite{nistGCM}. GCM is an instance
of \emph{authenticated encryption with associated data} (AEAD): in a single pass
it provides both confidentiality (an adversary learns nothing about the
plaintext) and integrity/authenticity (any tampering with the ciphertext is
detected on decryption, which fails rather than returning corrupted plaintext).

AEAD is the correct primitive here for two reasons. First, a released payload may
have travelled through untrusted storage and untrusted networks before a
journalist decrypts it; the authentication tag guarantees they are decrypting the
exact bytes the source sealed, not a forgery. Second, GCM requires a unique
\emph{initialization vector} (IV, also called a nonce) per encryption under a
given key; reusing an IV in GCM is catastrophic, so the implementation generates
a fresh 96-bit random IV for every payload (Section~\ref{sec:impl-encryption}).
A 96-bit IV is the size recommended by NIST for GCM because it is used directly
without additional processing.

A symmetric scheme alone, however, only relocates the problem: whoever holds the
AES key can read the data. The architectural question is therefore not ``how do
we encrypt'' but ``who is allowed to hold the key, and under what condition''.
The answer---threshold cryptography---requires the next two sections.

\section{Elliptic-Curve Cryptography}
\label{sec:bg-ecc}

\emph{Public-key} (asymmetric) cryptography uses a key pair: a private key kept
secret and a public key shared freely. \emph{Elliptic-curve cryptography}
(ECC)~\cite{koblitz1987, miller1986} achieves the security of older schemes such
as RSA at far smaller key sizes by working in the group of points on an elliptic
curve over a finite field. Ethereum and Bitcoin both use the curve
\texttt{secp256k1}~\cite{sec2}, and this work uses it as well for stealth
addresses.

Two operations matter for the present work. \emph{Scalar multiplication} takes a
private scalar $k$ and the curve's fixed generator point $G$ and computes the
public key $P = kG$; recovering $k$ from $P$ is the elliptic-curve
discrete-logarithm problem and is believed infeasible. \emph{Elliptic-curve
Diffie--Hellman} (ECDH) lets two parties who each hold a key pair derive a shared
secret without transmitting it: given private $a$ and the other party's public
$B = bG$, party $A$ computes $aB = abG$, and party $B$ symmetrically computes
$bA = baG = abG$; the two results are equal. Stealth addresses
(Section~\ref{sec:bg-stealth}) are built directly on these two operations: a
one-time payment address is derived by offsetting a recipient's public key by a
scalar that only the sender and recipient can compute via ECDH.

\section{Threshold Cryptography and Multi-Party Computation}
\label{sec:bg-threshold}

The core of the project's trust model is \emph{threshold cryptography}, a branch
of \emph{secure multi-party computation} (MPC) in which a cryptographic
capability is distributed among $n$ parties such that any $t$ of them can
exercise it together, but any fewer than $t$ learn nothing.

The foundational building block is \emph{Shamir's Secret
Sharing}~\cite{shamir1979}. To share a secret $s$ among $n$ parties with
threshold $t$, one constructs a random polynomial of degree $t-1$ whose constant
term is $s$, and gives each party the value of the polynomial at a distinct
point. Any $t$ parties hold $t$ points, which uniquely determine the degree-$t-1$
polynomial (by Lagrange interpolation) and therefore reveal $s$; any $t-1$
parties hold $t-1$ points, through which infinitely many degree-$t-1$ polynomials
pass, so every value of $s$ remains equally likely. The scheme is
\emph{information-theoretically} secure below the threshold: fewer than $t$
shares give \emph{zero} information, independent of the adversary's computing
power.

A \emph{threshold cryptosystem}~\cite{desmedt1989, nistThreshold} applies this
idea so that the secret being shared is a \emph{private key} which is never
reconstructed in one place. Through \emph{distributed key generation} (DKG), the
$n$ parties jointly produce a public key while each holds only a share of the
corresponding private key; no party ever sees the whole private key, not even at
creation. To decrypt or sign, $t$ parties each compute a partial result using
their share, and the partials are combined into the final result---again without
ever assembling the key. A \emph{threshold signature scheme} (TSS) does this for
signatures; threshold decryption does it for decryption.

The security implications are exactly what a coercion-resistant system needs.
The private capability has no single physical location to seize, subpoena, or
attack. An adversary must simultaneously compromise at least $t$ of the $n$
parties---ideally operated by different entities in different
jurisdictions---to break confidentiality, and must take more than $n - t$ offline
to prevent the operation, so the system is both confidential and fault-tolerant
by construction. Critically, the strength of these guarantees comes from a
mathematical assumption about how many independent parties an adversary can
corrupt, not from a physical assumption about the impenetrability of a single
device. This distinction is the crux of the argument in
Section~\ref{sec:bg-tee} and Chapter~\ref{ch:evaluation}.

Schemes based on the Boneh--Lynn--Shacham (BLS) signature~\cite{boneh2001bls}
support \emph{non-interactive} threshold operations, in which each party can
produce its partial result independently and the partials are aggregated without
a multi-round protocol. This efficiency is what makes a network like Lit
Protocol practical as application infrastructure rather than a laboratory
construction.

\section{Zero-Knowledge Proofs and zkTLS}
\label{sec:bg-zk}

A \emph{zero-knowledge proof} (ZKP)~\cite{goldwasser1989} allows a \emph{prover}
to convince a \emph{verifier} that a statement is true while revealing nothing
beyond the truth of the statement itself. A ZKP must be \emph{complete} (a true
statement can always be proven), \emph{sound} (a false statement cannot be proven
except with negligible probability), and \emph{zero-knowledge} (the verifier
learns nothing other than validity). Modern non-interactive succinct proof
systems make the proof small and quick to verify, which is what allows a proof to
be checked inside a smart contract.

This thesis applies zero-knowledge proofs to the credibility problem through
\emph{zkTLS} (zero-knowledge transport-layer security). Ordinary TLS---the
\texttt{https} that secures the web---encrypts traffic between a client and a
server so that no eavesdropper can read it, but it produces no transferable
evidence of \emph{what} was exchanged: a screenshot of a corporate portal proves
nothing, because it can be trivially forged. zkTLS closes this gap. Building on
research such as DECO~\cite{zhang2020deco}, zkTLS protocols let a client come away
from a genuine TLS session with a cryptographic proof about the session's
contents---``the server at this domain returned a record stating my employment
status is active''---without revealing the rest of the session (their username,
password, session token, or any other field they choose to redact).

The \emph{Reclaim Protocol}~\cite{reclaim} implements this with a
\emph{witness} (an HTTPS proxy) that observes the encrypted handshake and
attests to its authenticity, while the client decrypts the response locally,
redacts everything except the target claim, and produces a zero-knowledge proof
binding the redacted claim to the authenticated session. The result separates the
\emph{fact} being proven from the \emph{identity} of the person proving it---
exactly the property a credible-but-anonymous source requires.
Section~\ref{sec:design-identity} describes how \zkw anchors such proofs
on-chain. A complementary technology, \emph{ZK-Email}~\cite{zkemail}, proves
properties of emails by verifying their DKIM signatures in zero knowledge; it is
discussed as future work in Chapter~\ref{ch:conclusions}.

\section{Stealth Addresses (ERC-5564)}
\label{sec:bg-stealth}

A naive on-chain payment to a whistleblower is self-defeating: the recipient's
address is public, so paying it links the payment---and everyone who can see the
chain---to the source. \emph{Stealth addresses}~\cite{erc5564} solve this using
the ECDH machinery of Section~\ref{sec:bg-ecc}.

The recipient publishes a \emph{stealth meta-address} consisting of two public
keys: a \emph{spending} key and a \emph{viewing} key. To pay them, a sender
generates a fresh \emph{ephemeral} key pair, performs ECDH between the ephemeral
private key and the recipient's viewing public key to obtain a shared secret,
hashes that secret to a scalar, and offsets the recipient's spending public key
by that scalar times $G$ to produce a brand-new, one-time \emph{stealth address}.
The sender publishes only the ephemeral public key (an ``announcement''). The
recipient, scanning announcements, performs the symmetric ECDH with their viewing
private key, recomputes the same scalar, and recognises payments meant for them;
because they hold the spending private key, they---and only they---can derive the
private key controlling the stealth address and move the funds. A \emph{view tag}
(a single byte of the shared secret) lets the recipient skip most announcements
cheaply during scanning.

To any external observer, a stealth payment is a transfer to a fresh,
never-before-seen address with no visible connection to the recipient's published
identity. This is exactly the unlinkability the marketplace
(Section~\ref{sec:design-marketplace}) needs to reward sources safely.

\section{Permanent Decentralised Storage}
\label{sec:bg-storage}

Encrypted evidence must survive even a determined effort to erase it, and it must
remain available long after the source has gone silent. Conventional cloud
storage fails both tests: it can be deleted by the provider, by a court order, or
by the source under duress. \emph{Arweave}~\cite{arweave} addresses this with a
\emph{pay-once, store-forever} model: a single up-front payment funds an
endowment that incentivises a decentralised network of miners to replicate the
data indefinitely. Once a payload is written to Arweave's ``permaweb'', it cannot
be unilaterally deleted by anyone---including its author---which is precisely the
property a dead man's switch requires, since the source must not be able to
retract the evidence even when coerced.

A related system, the \emph{InterPlanetary File System}
(IPFS)~\cite{benet2014ipfs}, provides content-addressed storage where data is
referenced by the hash of its contents; it is well-suited to mutable metadata but
does not by itself guarantee permanence. The implementation
(Section~\ref{sec:impl-storage}) stores the encrypted vault payload on Arweave,
reached through a modern upload service that removes the historical friction of
funding storage transactions.

\section{Trusted Execution Environments and Their Limitations}
\label{sec:bg-tee}

The final piece of background is the technology this thesis deliberately moves
\emph{away} from. A \emph{trusted execution environment} (TEE) is a hardware
feature---Intel's \emph{Software Guard Extensions} (SGX) being the canonical
example~\cite{costan2016sgx}---that creates an \emph{enclave}: a region of
encrypted memory isolated from the operating system, the hypervisor, and even a
malicious machine operator. Code running in an enclave can hold secrets, such as
decryption keys, and prove to a remote party through \emph{attestation} that it
is running unmodified inside genuine hardware. On its face this is attractive for
key custody: the host runs the code but cannot see the keys.

The earlier iteration of this project relied on exactly this, sealing decryption
keys in an SGX enclave accessed through the iExec network. For the threat model of
high-value whistleblowing, however, a TEE has three fundamental weaknesses.

\textbf{Side-channel and transient-execution attacks.} The isolation guarantee of
SGX has been repeatedly broken in practice. Attacks such as
Foreshadow~\cite{vanbulck2018foreshadow} and SGAxe/CacheOut~\cite{vanschaik2020sgaxe}
exploit speculative execution and cache timing to extract sealed secrets,
including the attestation keys that underpin the whole trust chain. An adversary
with physical or privileged access to the host---a plausible position for a
nation-state acting against a cloud provider under legal compulsion---is exactly
the adversary best placed to mount such attacks.

\textbf{A centralised hardware root of trust.} SGX attestation ultimately
requires trusting the hardware manufacturer to have generated and managed the
attestation keys correctly and not to have been compelled to produce a backdoor.
This reintroduces, in silicon, the very single point of trust that
decentralisation is meant to remove: the security of the whole system collapses
if one vendor's master keys are compromised or coerced.

\textbf{Liveness and centralisation of operation.} TEE-based services tend to run
on specific, specialised hardware operated by a small number of providers. If the
particular nodes holding a source's key are taken offline or censored, the dead
man's switch can simply fail to fire---an unacceptable failure mode for a
mechanism whose entire value proposition is that it is unstoppable.

The conclusion drawn in this thesis, and developed in Chapter~\ref{ch:evaluation},
is that a TEE relocates trust from software to hardware but does not
\emph{eliminate} a single point of failure---it merely changes its nature, from
one that can be hacked to one that can be both hacked and legally compelled.
Threshold cryptography, by contrast, replaces the single point of failure with a
distributed-trust assumption that an adversary must defeat in many places at once.
It is on this basis that the system replaces the SGX enclave with the Lit
Protocol's threshold network. The remainder of the thesis develops that design.
It is worth noting, and Chapter~\ref{ch:evaluation} returns to this, that parts
of the wider Lit ecosystem have since reintroduced TEE-based execution; the design
presented here consciously stays on the multi-party-computation path.
